\documentclass[12pt,a4paper]{report}

\usepackage{graphicx}
\usepackage{hyperref}
\usepackage{setspace}
\usepackage{amsmath}
\usepackage{amssymb}

\usepackage[italian]{babel}
\usepackage[italian]{cleveref}
\usepackage{pdfpages}

\hyphenation{sil-la-ba-zio-ne pa-ren-te-si}
\textwidth=450pt\oddsidemargin=0pt

\begin{document}

\title{
  Corso di Virtualizzazione e integrazione di Sistemi
  \\ Sviluppo servizio cloud
}
\author{
  Alessandro Rebosio \\
  Matricola: 1130557
}
\date{\today}

\maketitle

\pagenumbering{roman}
\setstretch{1.25}

\chapter*{Introduzione}
\label{chap:introduzione}
\addcontentsline{toc}{chapter}{Introduzione}
L'evoluzione del cloud computing ha trasformato radicalmente il modo
in cui gestiamo, archiviamo e iniziamo le informazioni, rendendo
l'accessibilità ubiqua un requisito fondamentale per le moderne
applicazioni. In questo contesto, il presente progetto,
sviluppato nell'ambito del corso di Virtualizzazione e Integrazione
di Sistemi, si pone l'obiettivo di progettare e implementare un
servizio cloud di archiviazione e gestione file scalabile e
containerizzato. Il sistema offre un'esperienza utente
reattiva e moderna, ispirata ai principali servizi di cloud storage,
permettendo la gestione gerarchica di file e cartelle, il caricamento
drag-and-drop e la gestione del cestino in modo sicuro.

L'architettura è stata studiata per enfatizzare la separazione delle
responsabilità e la portabilità: il sistema si avvale di Docker per
l'isolamento dei servizi, Nginx come reverse proxy per la gestione
sicura del traffico, Next.js per la logica full-stack
dell'applicazione e Supabase per la persistenza dei dati e la
gestione dell'autenticazione.

L'elaborato descrive le fasi di analisi, progettazione e
implementazione del servizio, illustrando i requisiti funzionali, le
scelte architetturali adottate e i vantaggi derivanti dall'impiego
della containerizzazione nello sviluppo e nel deployment di
applicazioni distribuite.

\setstretch{1}
\tableofcontents

\cleardoublepage
\pagenumbering{arabic}
\setcounter{page}{1}
\setstretch{1.25}

% CHAPTER 1
\chapter{Obiettivi e Requisiti di Progetto}
\label{cap:obiettivi}

\section{Requisiti Funzionali}
Il progetto è stato concepito per offrire un servizio di cloud
storage completo e versatile. Di seguito vengono dettagliate le
funzionalità principali del sistema:

\begin{itemize}
  \item \textbf{Autenticazione e Gestione Utenti:} Sistema di
    registrazione e autenticazione robusto, gestito tramite i servizi
    di Supabase Auth, comprensivo di supporto per il recupero e il
    ripristino delle password.
  \item \textbf{Gestione dei File (Upload e Download):} Possibilità
    per gli utenti autorizzati di caricare file e intere directory
    (anche tramite interazione drag-and-drop) e di scaricare i propri
    contenuti in modo sicuro attraverso URL temporanei presignati.
  \item \textbf{Organizzazione e Navigazione:} Struttura gerarchica
    dei dati che permette la creazione di cartelle e sottocartelle,
    navigabili in modo intuitivo attraverso l'uso di breadcrumbs.
  \item \textbf{Meccanismo di Eliminazione (Cestino):}
    Implementazione di un sistema di \textit{soft-delete} che sposta
    i file eliminati in un'apposita sezione ``Cestino'',
    consentendone il recupero prima dell'eliminazione definitiva.
  \item \textbf{Gestione dei Ruoli e Permessi:} Gestione
    differenziata degli accessi tra utenti standard (responsabili
    unicamente del proprio spazio di archiviazione) e amministratori
    (abilitati a invitare nuovi utenti o a revocarne l'accesso).
\end{itemize}

\section{Requisiti Non Funzionali}
Per garantire l'affidabilità e le prestazioni del sistema, sono stati
definiti rigorosi requisiti non funzionali che guidano l'architettura
dell'applicazione:

\begin{itemize}
  \item \textbf{Sicurezza:} Gestione sicura delle sessioni utente e
    controllo d'accesso ai dati regolato dalle policy di Row Level
    Security (RLS) di PostgreSQL, che garantiscono l'isolamento delle
    risorse per ciascun utente.
  \item \textbf{Scalabilità e Deployment:} Architettura interamente
    strutturata in un ecosistema a container basato su Docker, utile
    a gestire in modo flessibile il deployment in base all'ambiente
    di esecuzione.
  \item \textbf{Affidabilità e Persistenza dei Dati:} Utilizzo di un
    database relazionale robusto per la consistenza dei metadati e di
    un'architettura di object storage resiliente per la conservazione
    dei file binari.
  \item \textbf{Usabilità e Interfaccia Utente:} Interfaccia grafica
    sviluppata con React e Tailwind CSS, completamente responsive per
    garantire un'esperienza d'uso ottimale su qualsiasi tipo di dispositivo.
  \item \textbf{Manutenibilità e Qualità del Codice:} Codice sorgente
    modulare e fortemente tipizzato grazie all'uso di TypeScript,
    finalizzato a facilitare l'estendibilità e il debugging del sistema.
\end{itemize}

\begin{figure}[htbp]
  \centering
  \includegraphics[width=0.75\textwidth]{./img/requirements.png}
  \caption{Panoramica sinottica dei requisiti funzionali e non
  funzionali del sistema.}
  \label{fig:requirements}
\end{figure}

\section{Vincoli Tecnologici}
Le scelte progettuali sono state fortemente indirizzate da specifici
vincoli tecnologici e obiettivi didattici, primo fra tutti l'adozione
della containerizzazione. Sul fronte dello stack applicativo, il
progetto richiede l'utilizzo del framework moderno Next.js per
gestire in un'unica base di codice sia l'interfaccia grafica React
sia le logiche di backend. Il sistema si appoggia inoltre alla suite
Supabase per fornire i servizi di database PostgreSQL, lo storage
compatibile S3 e il modulo di autenticazione.

Per garantire la massima flessibilità, il progetto prevede la
gestione di due ambienti distinti attraverso l'isolamento dei
processi, ma con architetture differenti per assecondare le
specifiche esigenze della fase del ciclo di vita del software:

\subsection{Ambiente di Sviluppo (Dev)}
L'ambiente di sviluppo locale è orchestrato tramite il file
\texttt{docker-compose.dev.yml}. Questa configurazione è ottimizzata
per agevolare le attività di programmazione e debugging, offrendo
flessibilità e strumenti di monitoraggio immediati per lo
sviluppatore durante la scrittura del codice.

In questo contesto, il flusso di lavoro è supportato dall'utilizzo
della \textit{Supabase CLI} (sviluppata in Node.js), la quale
semplifica la gestione dell'infrastruttura locale avviando e
coordinando automaticamente l'intero ecosistema di container Docker
richiesti da Supabase (come il database PostgreSQL, l'autenticazione
e lo storage). Ciò permette di replicare fedelmente i servizi cloud
direttamente sulla macchina locale di sviluppo.

\subsection{Ambiente di Produzione (Prd)} % TODO
L'ambiente di produzione è definito nel file
\texttt{docker-compose.yml}. Questa architettura impone criteri
rigorosi di sicurezza e stabilità; in particolare, prevede l'utilizzo
di Nginx configurato come reverse proxy, il cui compito è instradare
in sicurezza le richieste HTTP in ingresso isolando i servizi interni
dalla rete pubblica.

% CHAPTER 2
\chapter{Architettura del Sistema}
\label{cap:architettura}

\section{Architettura a Container e Ambienti Operativi}
L'architettura del sistema fa perno sulla containerizzazione per
garantire portabilità e flessibilità, adottando due configurazioni
distinte per gli ambienti di sviluppo e di produzione. Nell'ambiente
di sviluppo, configurato tramite \texttt{docker-compose.dev.yml},
l'attenzione è posta sulla rapidità di iterazione. Il container della
web-app avvia il server in modalità di sviluppo, permettendo
l'\textit{hot-reloading} del codice, mentre l'intera infrastruttura
di backend e database viene demandata all'esecuzione locale della
Supabase CLI, che emula l'ecosistema cloud direttamente sulla
macchina dello sviluppatore senza appesantire la configurazione
Docker principale del progetto.

L'ambiente di produzione, governato dal file
\texttt{docker-compose.yml}, implementa invece un'architettura
self-hosted autonoma e completa. In questo scenario, l'infrastruttura
si espande includendo non solo il container \texttt{nginx-proxy} come
punto di ingresso per instradare le richieste alla \texttt{web-app},
ma ospita direttamente tutti i servizi core della suite Supabase.
Vengono istanziati container dedicati per il database PostgreSQL, per
le API REST fornite da PostgREST, per la gestione dei file binari
tramite lo Storage e per il modulo Auth, realizzando un ecosistema
completamente indipendente e isolato all'interno della rete Docker.

\section{Comunicazione tra Container}
La comunicazione all'interno del sistema avviene sfruttando la rete
virtuale privata gestita dinamicamente da Docker. Nell'ambiente di
produzione, questo isolamento è fondamentale poiché solo il proxy
Nginx espone la porta HTTP (80) verso l'esterno, mentre tutti gli
altri container dialogano esclusivamente tramite hostname interni. Il
flusso delle comunicazioni inizia quando il browser dell'utente
finale invia una richiesta HTTP all'indirizzo pubblico. Nginx
intercetta la richiesta, vi inietta gli header informativi, come
l'indirizzo IP reale del client, e la inoltra alla web-app Next.js. A
sua volta, l'applicazione Next.js si interfaccia con i servizi
backend. In produzione, queste richieste vengono indirizzate verso i
container auto-ospitati: le interrogazioni relazionali passano
attraverso PostgREST per accedere a PostgreSQL, i caricamenti dei
file vengono gestiti dal container Storage e le validazioni dei token
dal servizio Auth, il tutto senza che alcun dato transiti sulla rete
dell'host fisico o su internet.

\section{Infrastruttura e Networking in Produzione}
L'inserimento di Nginx nel ruolo di reverse proxy nell'ambiente di
produzione riveste un'importance strategica. Nginx opera come una
barriera protettiva che disaccoppia la rete pubblica dalla rete
privata in cui operano i container. Oltre a garantire la sicurezza
nascondendo le reali porte e gli indirizzi dei server di backend, il
proxy gestisce in modo altamente efficiente i buffer per le risposte
e l'aggiornamento dei protocolli per i WebSocket. Questa netta
separazione delle responsabilità mitiga il rischio di esposizione
diretta dell'ambiente Node.js e dell'infrastruttura Supabase a
vulnerabilità esterne e predispone il sistema per un'agevole
scalabilità orizzontale tramite tecniche di load balancing.

\newpage
\section{Diagramma Architetturale}
A completamento della descrizione dettagliata dei singoli componenti
e dei flussi di rete, viene presentata in questa sezione una vista
d'insieme dell'infrastruttura di produzione.

La Figura~\ref{fig:architecture} illustra la topologia di
rete del sistema all'interno dell'ecosistema Docker. Il diagramma
evidenzia chiaramente la separazione dei compiti e il flusso dei
dati: l'utente finale interagisce esclusivamente con il punto di
ingresso pubblico (Nginx), il quale si occupa di smistare le
richieste di pagine Web verso il frontend Next.js e il traffico API
verso il gateway Kong. Quest'ultimo, cuore pulsante dello stack
auto-ospitato di Supabase, valida i token JWT e instrada le richieste
verso i microservizi interni (\textit{GoTrue, PostgREST e Storage}),
i quali condividono l'accesso sicuro e isolato al database
relazionale PostgreSQL regolato dalle policy RLS.

\begin{figure}[htbp]
  \centering
  \includegraphics[width=\textwidth]{./img/architecture.png}
  \caption{Diagramma architetturale dell'infrastruttura in ambiente
  di produzione.}
  \label{fig:architecture}
\end{figure}

% CHAPTER 3
\chapter{Containerizzazione e Deployment}
\label{cap:docker}

\section{Struttura dei Container}
% Descrivere:
% - immagini Docker utilizzate
% - configurazione dei servizi
% - dipendenze tra container

\section{Orchestrazione con Docker Compose}
% Spiegare:
% - come vengono avviati i servizi
% - definizione delle reti
% - collegamento tra container

\section{Gestione dei Volumi}
% Descrivere:
% - persistenza dati PostgreSQL
% - persistenza storage (filesystem / MinIO)
% - differenza tra volumi e storage effimero

\section{Vantaggi della Containerizzazione}
% Analizzare:
% - isolamento dei servizi
% - portabilità
% - facilità di deployment
% - confronto con architettura monolitica

% CHAPTER 4
\chapter{Protocolli e Tecnologie Utilizzate}
\label{cap:protocolli}

\section{Comunicazione HTTP e REST}
% Descrivere:
% - uso di HTTP per le API
% - architettura REST
% - metodi principali (GET, POST)
% - upload file tramite multipart/form-data

\section{Comunicazione tra Servizi}
% Descrivere:
% - comunicazione backend → database
% - comunicazione backend → storage
% - uso di API S3 nel caso MinIO

\section{Persistenza dei Dati}
% Confrontare:
% - filesystem (struttura gerarchica)
% - object storage (bucket e oggetti)
% Evidenziare differenze in accesso e scalabilità.

% CHAPTER 5
\chapter{Gestione Utenti e Autenticazione}
\label{cap:autenticazione}

\section{Modello di Autenticazione JWT}
% Descrivere:
% - generazione token
% - validazione
% - sessione stateless
% Inserire diagramma di sequenza del login.

\section{Sicurezza delle Credenziali}
% Spiegare:
% - hashing password (bcrypt)
% - autorizzazione accessi
% - protezione delle risorse

% CHAPTER 6
\chapter{Possibili Estensioni e Visione Futura}
\label{cap:estensioni}

\section{Gestione Permessi}
% Descrivere:
% - accessi ai file
% - condivisione tra utenti

\section{Scalabilità e Cloud}
% Descrivere:
% - migratione verso cloud reale
% - vantaggi e criticità

% CONCLUSIONI
\chapter*{Conclusioni}
\addcontentsline{toc}{chapter}{Conclusioni}
% Riassumere:
% - obiettivi raggiunti
% - differenze tra filesystem e object storage
% - vantaggi della containerizzazione
% - competenze acquisite (sistemi distribuiti, sicurezza, architettura)

\end{document}
